% !TEX program = xelatex
% 공통수학2 · Ⅱ. 집합과 명제 · 2. 명제 · 03 명제의 증명 · 1/2차시
% 학생용 활동지 (답안 없음)
\documentclass[11pt,a4paper]{article}
\usepackage{kotex}
\usepackage{iftex}
\ifPDFTeX\else
  \setmainhangulfont{Noto Serif CJK KR}
  \setsanshangulfont{Noto Sans CJK KR}
\fi
\usepackage[margin=2.1cm]{geometry}
\usepackage{parskip}
\usepackage{enumitem}
\usepackage{array}
\usepackage{tabularx}
\usepackage{xcolor}
\usepackage{amssymb}
\usepackage{tikz}
\usepackage[most]{tcolorbox}
\usepackage{fancyhdr}
\usepackage{titlesec}

\definecolor{goldc}{HTML}{B07C22}
\definecolor{goldstrong}{HTML}{8A5F16}
\definecolor{indigoc}{HTML}{2B3B57}
\definecolor{indigosoft}{HTML}{6E82A6}
\definecolor{linec}{HTML}{D6DACB}
\definecolor{surface2}{HTML}{E4E8DC}
\definecolor{writeline}{HTML}{C7CCBA}
\definecolor{inksoft}{HTML}{52604F}

\titleformat{\section}{\normalfont\sffamily\bfseries\large\color{indigoc}}{\thesection}{0.6em}{}
\titlespacing{\section}{0pt}{14pt}{6pt}
\newtcolorbox{goalbox}{enhanced, breakable, colback=white, colframe=goldc, boxrule=0.5pt, leftrule=3pt, arc=2pt, left=10pt, right=10pt, top=8pt, bottom=8pt}
\newtcolorbox{sheetbox}{enhanced, breakable, colback=white, colframe=linec, boxrule=0.6pt, arc=3pt, left=11pt, right=11pt, top=9pt, bottom=9pt}
\newcommand{\writespace}[1]{%
  \par\nobreak\vspace{3pt}
  \foreach \n in {1,...,#1} {%
    \noindent\textcolor{writeline}{\rule{\linewidth}{0.4pt}}\par\vspace{0.62cm}%
  }
}
\newcommand{\fillblank}[1]{\makebox[#1]{\hrulefill}}
\pagestyle{fancy}\fancyhf{}\renewcommand{\headrulewidth}{0pt}
\fancyfoot[L]{\footnotesize\color{inksoft} 공통수학2 · 2026학년도 2학기 · 지평선고등학교}
\fancyfoot[R]{\footnotesize\color{inksoft} 다음 시간 --- 03 명제의 증명(2)}

\begin{document}

\noindent
\begin{minipage}[b]{0.62\linewidth}
  {\sffamily\footnotesize\color{indigosoft} 공통수학2 · Ⅱ. 집합과 명제 · 2. 명제}\\[2pt]
  {\sffamily\bfseries\Large 03 명제의 증명 \textperiodcentered\ 30차시}
\end{minipage}\hfill
\begin{minipage}[b]{0.36\linewidth}
  \raggedleft\small\color{inksoft}
  반\ \fillblank{1.6cm} \quad 번호\ \fillblank{1.6cm}\\[4pt]
  이름\ \fillblank{3.6cm}
\end{minipage}

\vspace{4pt}
\noindent{\color{goldc}\rule{\linewidth}{2.2pt}}
\vspace{10pt}

\begin{goalbox}
\textbf{오늘의 목표}\\
대우를 이용한 증명법과 귀류법을 이해하고, 관련된 명제를 증명할 수 있다.
\vspace{4pt}
\small\color{inksoft}
$\square$ 자신 있음 \quad $\square$ 조금 헷갈림 \quad $\square$ 처음 봄 \hfill (수업 전 나의 상태에 표시)
\end{goalbox}

\section{생각 열기 --- 놀이공원 할인}

\begin{sheetbox}
어느 놀이공원에서는 그 지역 주민 또는 교복을 입은 사람에게 입장료를 할인해 준다.

\begin{enumerate}[leftmargin=1.4em, itemsep=2pt]
  \item `교복을 입은 사람은 입장료를 할인받는다.'가 참인 문장인지 말해 보자.
  \item `입장료를 할인받지 못한 사람은 교복을 입고 있지 않다.'가 참인 문장인지 말해 보자.
\end{enumerate}
\writespace{3}
\end{sheetbox}

\section{개념 정리 --- 대우를 이용한 증명법과 귀류법}

\begin{sheetbox}
어떤 명제가 참이면 그 \fillblank{2.2cm}도 참이고, 대우가 참이면 처음의 명제도 참이다. 따라서 어떤 명제가 참임을 증명하는 대신 그 \fillblank{2.2cm}가 참임을 증명해도 된다.

\vspace{8pt}
\textbf{예제} --- `자연수 $n$에 대하여 $n^2$이 짝수이면 $n$도 짝수이다.'를 대우로 증명하시오.\\
증명: 주어진 명제의 대우는 `자연수 $n$에 대하여 $n$이 홀수이면 $n^2$도 홀수이다.'이다. $n$이 홀수이면 $n=2k-1$($k$는 자연수)로 나타낼 수 있다. 이때\\
\hspace*{2em}$n^2=(2k-1)^2=4k^2-4k+1=2(2k^2-2k+1)-1$\\
이므로 $n^2$도 홀수이다. 따라서 대우가 참이므로 주어진 명제도 참이다.

\vspace{8pt}
어떤 명제가 참임을 증명할 때, 그 명제의 부정이 참이라고 가정하면 이미 알려진 사실에 \fillblank{2.2cm}이 생긴다는 것을 보이는 방법을 \fillblank{2.6cm}이라고 한다.
\end{sheetbox}

\section{연습 문제}

\begin{sheetbox}
\textbf{1.} 대우를 이용하여 다음 명제가 참임을 증명하시오.\\
\hspace*{1.6em}`자연수 $n$에 대하여 $n^2$이 3의 배수이면 $n$도 3의 배수이다.'\\
\small(힌트: $n$이 3의 배수가 아니면 $n=3k-2$ 또는 $n=3k-1$로 나타낼 수 있다.)
\writespace{7}

\vspace{6pt}
\textbf{2.} 귀류법을 이용하여 다음 명제가 참임을 증명하시오.\\
\hspace*{1.6em}`$1+\sqrt2$는 유리수가 아니다.'
\writespace{6}
\end{sheetbox}

\section{사고의 가시화 --- Connect · Extend · Challenge}

\noindent{\footnotesize\color{inksoft} 오늘 배운 `대우를 이용한 증명법과 귀류법'을 떠올리며 아래 세 칸을 채워 보자.}\\[6pt]
\begin{tabularx}{\linewidth}{@{}XXX@{}}
\textbf{\footnotesize\color{indigoc} Connect --- 무엇과 연결되나?} &
\textbf{\footnotesize\color{indigoc} Extend --- 어디까지 확장됐나?} &
\textbf{\footnotesize\color{indigoc} Challenge --- 아직 궁금한 것은?} \\[4pt]
\end{tabularx}
\noindent
\begin{minipage}[t]{0.32\linewidth}\writespace{3}\end{minipage}\hfill
\begin{minipage}[t]{0.32\linewidth}\writespace{3}\end{minipage}\hfill
\begin{minipage}[t]{0.32\linewidth}\writespace{3}\end{minipage}

\section{마무리 성찰}

\begin{sheetbox}
\begin{minipage}[t]{0.48\linewidth}
{\footnotesize\color{inksoft} 오늘 배운 것을 한 문장으로}
\writespace{2}
\end{minipage}\hfill
\begin{minipage}[t]{0.48\linewidth}
{\footnotesize\color{inksoft} 감사 일기 / 유념 한 줄}
\writespace{2}
\end{minipage}
\end{sheetbox}

\end{document}
